Kerja sama dalam penelitian ini memiliki urgensi yang tinggi mengingat kebutuhan akan dataset BISINDO yang valid, beretika, dan merepresentasikan praktik bahasa isyarat yang digunakan oleh komunitas Tuli.
\begin{enumerate}[noitemsep]
    \item Proses pembentukan dataset bahasa isyarat tidak hanya memerlukan kemampuan teknis, tetapi juga pemahaman linguistik dan konteks sosial yang tidak dapat dipisahkan dari peran komunitas dan lembaga terkait.
    \item Bagi Pusat Bahasa Isyarat Indonesia (PUSBISINDO), kerja sama ini berpotensi memberikan kontribusi dalam bentuk dokumentasi terstruktur bahasa isyarat berbasis data \textit{landmark} yang lebih aman secara privasi. Keterlibatan PUSBISINDO dalam proses validasi dan pendampingan memungkinkan dataset yang dihasilkan tetap selaras dengan kaidah BISINDO serta nilai yang dijunjung oleh komunitas Tuli.
    \item Bagi komunitas Tuli, penelitian ini diharapkan dapat mendukung pengembangan teknologi bantu yang tidak menggantikan peran manusia, melainkan berfungsi sebagai sarana pendukung komunikasi dan pembelajaran. Pendekatan berbasis \textit{landmark} membantu mengurangi risiko eksploitasi visual sekaligus menjaga martabat partisipan dalam proses pengambilan data.
    \item Bagi pengembangan akademik, penelitian ini memberikan kontribusi pengetahuan terkait pemanfaatan model spasial dua dimensi dalam pengenalan bahasa isyarat berbasis koordinat gerakan. Hasil penelitian diharapkan dapat menjadi referensi bagi kajian lanjutan yang mengeksplorasi pendekatan ringan dan efisien dalam pemodelan gerakan manusia, khususnya pada konteks bahasa isyarat Indonesia.
\end{enumerate}
